\documentclass[10pt]{article} 
\renewcommand{\thempfootnote}{\arabic{mpfootnote}}

\usepackage[utf8]{inputenc}

\usepackage{multicol}
\setlength{\columnsep}{1cm}

\usepackage{abc}
\usepackage{graphicx}
\usepackage{musicography}

\usepackage[margin=0.5in]{geometry} 
\geometry{letterpaper}

\usepackage{fancyhdr} 
\pagestyle{fancy} % options: empty , plain , fancy
\renewcommand{\headrulewidth}{0pt} % customise the layout...
\lhead{}\chead{}\rhead{}
\lfoot{}\cfoot{\thepage}\rfoot{}

\begin{document}

\thispagestyle{empty}

\begin{abc}[name=thereunion,options={-O= +c -F single.fmt}]
X:1
T:The Reunion
C:Vivian Williams
M:C
Z:abc-transcription: Josh Larios <jdlarios@gmail.com>, 2021.08.02
K:D
AG | "D"F2 !uppermordent!ED d3 c | "G"BA Bd "D"A2 DE | FA "G"Bd "D"AF ED | "Em"B,2 "E"E2 "A7"E2 AG | 
"D"F2 EF DF Ad | "G"!uppermordent!BA B<d "A7"A3 G | "D"Fd "G"BG "A7"AF EG | "D"F2 "G"D2 "D"D2 :|
|: Bc | "D"d3 f "A"ec B<A | "G"BG DB "A"A3 "A7"c | "D"d2 cd "G"BA "D"df | "E7"eB ^Ge "A"cA "A7"fg | 
"D"af dA "G"G2 B<d | "D"AF Ad "Em"G3 "A7"A | "D"Fd "G"BG "A7"AF EG | "D"F2 "G"D2 "D"D2 :|\end{abc}
\vskip 0.5in
\noindent\makebox[\textwidth][c]{%
  \begin{minipage}{7.0in}
   \setlength{\parindent}{1.5em}

   We are so honored and privileged to feature one of the Pacific
   Northwest's own compositions for this month's WOTFA Tune of the
   Month. \textbf{The Reunion} is an air composed by Vivian Williams and was
   originally published in \textbf{151 Brand New Old Time Fiddle Tunes, Volume
   1}, edited by Vivian T. Williams, Voyager Publications, 1983. The
   opening 6 measures were used in the \textbf{Missing Fiddle Camp 2020}
   t-shirt, which was a successful fundraiser for WOTFA during the
   global pandemic.

   We were fortunate enough to talk to Vivian about her background and
   music composition philosophy.

   \textit{From Kitty -- It is my understanding that you have an extensive
   classical music background. How was it that you slipped away from
   classical music and towards folk music?}

   I took classical music lessons from the age of 9 through college,
   most liking baroque music and chamber music. Got interested in
   traditional folk music in college, but I didn’t know anything about
   fiddling — I was thinking guitars, banjos, Pete Seeger kind of
   stuff. When Phil and I got married in 1959 and moved to Seattle, we
   got involved with the folk music scene, and didn’t have any
   contacts with the classical music scene. So we went into it
   wholeheartedly. I heard some bluegrass fiddling on records, but
   didn’t particularly like it -- it sounded to harsh and dissonant to
   my ears. But then I heard a record of Mike Seeger playing ``Old
   Molly Hare'' -- a totally sweet but simple ``kindergarten level'' tune,
   and liked the sound of it. I figured ``I could do that!!'' -- and the
   rest is history!

   \textit{From Katie - I'm curious about inspiration. How did the tune come to exist?}

   Frank Ferrel introduced me to the Cape Breton fiddling of Winston
   Fitzgerald and Jerry Holland, and I really liked the way that style
   of fiddling combined classical technique with gutsy dance
   rhythms. And I enjoyed their medleys which would start with an air,
   then a strathspey, and end with a reel. So ``The Reunion'' was
   written as the first tune of that kind of a medley.

   \textit{From Josh - I'd be interested to know what Cape Breton /
     Scottish fiddlers you were influenced by when you wrote this
     tune, and if there's anyone you'd recommend people listen to
     now.}

   Winston Fitzgerald, Jerry Holland (whom I met at the Festival of
   American Fiddle Tunes in Port Townsend), Carl MacKenzie, and a
   woman fiddler whose name I've forgotten, but she played airs
   beautifully. There are plenty of others worth listening to, but
   these were my favorites at the time.

   \textit{From Kitty - The EM – AM progression in measure 13 is quite
     lovely and gives the final cadence an even stronger catharsis. Do
     you have an opinion on which chord progressions work better with
     certain tune types? Or should melody or rhythm rule all?}

   Gee, I dunno, in a tune like this, rhythm doesn't count for a lot,
   but in a dance tune it's primary. In an air, melody is
   important. Chords are there to enhance the melody, or to add
   interest and excitement to any kind of tune. But although chords
   can be important, I think they are secondary, in this style of
   music (compared with jazz or bluegrass, where the chord progression
   is basic to the tune).
   
   
 \end{minipage}
}


\end{document}
